\pagestyle{empty}
\tight
\begin{tblr}{width=\textwidth,colspec={|Q[c,m,1.9cm]|X[l,m]|},hlines,vlines,hspan=minimal,row{1}={bg=headblue},rowsep=1pt,colsep=3pt}
\SetCell{c,m}\hfil\logo\hfil & \textbf{POLITEKNIK NEGERI MALANG}\par\textbf{JURUSAN TEKNOLOGI INFORMASI}\par\textbf{PROGRAM STUDI: D4 TEKNIK INFORMATIKA} \\
\SetCell[c=2]{c} \textbf{RENCANA PEMBELAJARAN SEMESTER (RPS)} & \\
\end{tblr}
\begin{longtblr}{width=\textwidth,colspec={|Q[l,m,1.9cm]|X[0.85,l,m]|X[1.45,l,m]|X[0.75,l,m]|X[1.15,l,m]|},hlines,vlines,hspan=minimal,rowsep=1pt,colsep=2pt,row{1,3}={bg=headgray,font=\bfseries}}
MATA KULIAH & KODE & BOBOT (SKS)/jam & SEMESTER & TGL. PENYUSUNAN \\
Pengolahan Citra dan Visi Komputer & RTI255007 & 3 SKS/6 jam & 5 & 1 Juli 2025 \\
OTORISASI & Dosen Pengembang RPS & Koordinator RMK & \SetCell[c=2]{l} Ka PRODI &  \\
 & Mamluatul Hani'ah, S.Kom., M.Kom. & Mamluatul Hani'ah, S.Kom., M.Kom. & \SetCell[c=2]{l}{\vspace{8mm}\par Dr. Ely Setyo Astuti, S.T., M.T.} &  \\
\SetCell[r=13]{l} Capaian Pembelajaran (CP) & \SetCell[c=4]{l,bg=headgray,font=\bfseries} Capaian Pembelajaran Lulusan Program Studi (CPL-Prodi) &  &  &  \\
 & CPL07 & \SetCell[c=3]{l} Mampu mengembangkan sistem komputasi cerdas dan melakukan analisis data berbasis kecerdasan buatan untuk pengambilan keputusan. &  &  \\
 & CPL10 & \SetCell[c=3]{l} Mampu menganalisis persoalan kompleks dengan wawasan multidisiplin untuk menghasilkan solusi di bidang informatika dan ilmu komputer. &  &  \\
 & \SetCell[c=4]{l,bg=headgray,font=\bfseries} Capaian Pembelajaran mata kuliah (CPL-MK) &  &  &  \\
 & CPMK0705 & \SetCell[c=3]{l} Mahasiswa mampu melakukan pengolahan citra digital dan visualisasi untuk menghasilkan representasi yang bermakna. &  &  \\
 & CPMK0706 & \SetCell[c=3]{l} Mahasiswa mampu mengembangkan sistem visi komputer berbasis deep learning untuk deteksi, segmentasi, dan pengenalan objek. &  &  \\
 & CPMK1004 & \SetCell[c=3]{l} Mahasiswa mampu mengevaluasi kualitas hasil pengolahan citra dan representasi visual untuk aplikasi nyata. &  &  \\
 & CPMK1005 & \SetCell[c=3]{l} Mahasiswa mampu mengevaluasi kinerja sistem visi komputer berdasarkan metrik kualitas yang sesuai konteks. &  &  \\
 & \SetCell[c=4]{l,bg=headgray,font=\bfseries} Kemampuan akhir tiap tahapan belajar (Sub-CPMK) &  &  &  \\
 & SCPMK0705-04301 & \SetCell[c=3]{l} Mahasiswa mampu mengimplementasikan teknik pengolahan citra (filtering, segmentasi, transformasi) menggunakan OpenCV. &  &  \\
 & SCPMK0706-04302 & \SetCell[c=3]{l} Mahasiswa mampu membangun sistem visi komputer berbasis CNN untuk klasifikasi, deteksi objek, dan segmentasi. &  &  \\
 & SCPMK1004-04303 & \SetCell[c=3]{l} Mahasiswa mampu mengevaluasi kualitas citra dan hasil pemrosesan visual menggunakan metrik kuantitatif. &  &  \\
 & SCPMK1005-04304 & \SetCell[c=3]{l} Mahasiswa mampu mengevaluasi kinerja model visi komputer menggunakan metrik mAP, IoU, dan accuracy. &  &  \\
\SetCell[r=2]{l} Penilaian & \SetCell[c=4]{l} *Nilai CPMK didapat dari agregasi nilai SUB-CPMK yang dibebankan &  &  &  \\
 & \SetCell[c=4]{l}{\tiny\begin{tblr}{width=\linewidth,colspec={|Q[c,m,0.1\linewidth]|Q[c,m,0.16\linewidth]|X[l,m]|X[l,m]|X[l,m]|X[l,m]|X[l,m]|X[l,m]|Q[c,m,0.08\linewidth]|},hlines,vlines,hspan=minimal,rowsep=1pt,colsep=0.7pt,row{1,2}={bg=headgray,font=\bfseries}}
Kode CPMK & Kode SCPMK & \SetCell[c=6]{c} Bobot per Bentuk Penilaian & & & & & & Total Bobot \\
 & & Tugas & Kuis 1 & UTS & Kuis 2 & PBL & UAS & \\
CPMK0705 & SCPMK0705-04301 & 12\%: filtering dan segmentasi & 5\%: pengolahan citra & 8\%: implementasi citra & 0\% & 0\% & 0\% & 25\% \\
CPMK0706 & SCPMK0706-04302 & 8\%: CNN awal & 0\% & 0\% & 0\% & 13\%: sistem visi & 11\%: evaluasi sistem & 32\% \\
CPMK1004 & SCPMK1004-04303 & 8\%: evaluasi kualitas citra & 0\% & 7\%: evaluasi citra & 0\% & 0\% & 0\% & 15\% \\
CPMK1005 & SCPMK1005-04304 & 0\% & 0\% & 0\% & 5\%: evaluasi model & 12\%: evaluasi sistem & 11\%: evaluasi sistem & 28\% \\
\SetCell[c=8]{r}\textbf{Total Per Penilaian} & & & & & & & & \textbf{100\%} \\
\end{tblr}} &  &  &  \\
Diskripsi Singkat Mata Kuliah & \SetCell[c=4]{l} Mata kuliah Pengolahan Citra dan Visi Komputer membangun kemampuan memproses citra digital dengan OpenCV dan mengembangkan sistem visi komputer berbasis CNN untuk klasifikasi, deteksi objek, dan segmentasi pada aplikasi nyata. &  &  &  \\
Bahan Kajian / Materi Pembelajaran & \SetCell[c=4]{l}\begin{enumerate}\item Pengolahan citra digital: representasi, operasi pixel, histogram, dan color space\item Filtering spasial, deteksi tepi, dan segmentasi citra\item Transformasi morfologi dan evaluasi kualitas citra\item CNN: arsitektur, klasifikasi, deteksi objek, dan segmentasi gambar\item Transfer learning, augmentasi data, dan evaluasi model visi komputer\end{enumerate} &  &  &  \\
\SetCell[r=2]{l} Pustaka & \SetCell[c=4]{l} \textbf{Utama:}\begin{enumerate}\item Gonzalez, R. 2018. Digital Image Processing (4th ed.). Pearson.\item Goodfellow, I. 2016. Deep Learning. MIT Press.\item Dokumentasi OpenCV dan PyTorch torchvision.\end{enumerate} &  &  &  \\
 & \SetCell[c=4]{l} \textbf{Pendukung:} Materi LMS, dokumentasi resmi perangkat lunak, dan jobsheet praktikum. &  &  &  \\
Media Pembelajaran & \SetCell[c=2]{l} \textbf{Software:} Python, OpenCV, PyTorch, Jupyter Notebook, LMS. &  & \SetCell[c=2]{l} \textbf{Hardware:} Komputer/laptop, GPU (jika tersedia), proyektor. &  \\
Nama Dosen Pengampu & \SetCell[c=4]{l} Tim Pengajar Program Studi D4 TI &  &  &  \\
Matakuliah Syarat & \SetCell[c=4]{l} RTI255004 Pembelajaran Mesin &  &  &  \\
\end{longtblr}
\begin{longtblr}{width=\textwidth,colspec={|Q[c,m,0.06\textwidth]|X[1.35,l,m]|X[1.08,l,m]|X[1.16,l,m]|X[0.7,l,m]|X[1.24,l,m]|X[1.06,l,m]|X[0.98,l,m]|Q[c,m,0.05\textwidth]|},hlines,vlines,hspan=minimal,rowhead=2,row{1,2}={bg=headgreen,font=\bfseries},rowsep=1pt,colsep=1.5pt}
Mgg.\par Ke & Kemampuan Akhir Yang Direncanakan (Sub-CP-MK) & Bahan kajian (Materi Pembelajaran) & Bentuk dan Metode Pembelajaran & Estimasi Waktu & Pengalaman Belajar Mahasiswa & Kriteria \& Bentuk Penilaian & Indikator Penilaian & Bobot Penilaian (\%) \\
(1) & (2) & (3) & (4) & (5) & (6) & (7) & (8) & (9) \\
1 & SCPMK0705-04301 & Pengantar pengolahan citra: representasi dan format. & Modalitas: Blended Learning. Bentuk: Luring/Daring. Metode: Case Method, praktik implementasi, eksplorasi dataset, dan peer review. & 1 x 3 x 50' tatap muka; 1 x 3 x 50' tugas/praktik mandiri. & Mahasiswa mengerjakan aktivitas terarah untuk pengantar pengolahan citra: representasi dan format dan menunjukkan evidence hasil belajar. & Bentuk penilaian: Pengantar Pengolahan Citra: Representasi dan Format. Mengacu rubrik RTM. & Ketepatan konsep; kualitas implementasi; validasi dan dokumentasi. & 2 \\
2 & SCPMK0705-04301 & Operasi pixel, histogram, dan color space. & Modalitas: Blended Learning. Bentuk: Luring/Daring. Metode: Case Method, praktik implementasi, eksplorasi dataset, dan peer review. & 1 x 3 x 50' tatap muka; 1 x 3 x 50' tugas/praktik mandiri. & Mahasiswa mengerjakan aktivitas terarah untuk operasi pixel, histogram, dan color space dan menunjukkan evidence hasil belajar. & Bentuk penilaian: Operasi Pixel, Histogram, dan Color Space. Mengacu rubrik RTM. & Ketepatan konsep; kualitas implementasi; validasi dan dokumentasi. & 3 \\
3 & SCPMK0705-04301 & Filtering spasial: konvolusi, Gaussian blur, sharpening. & Modalitas: Blended Learning. Bentuk: Luring/Daring. Metode: Case Method, praktik implementasi, eksplorasi dataset, dan peer review. & 1 x 3 x 50' tatap muka; 1 x 3 x 50' tugas/praktik mandiri. & Mahasiswa mengerjakan aktivitas terarah untuk filtering spasial: konvolusi, Gaussian blur, sharpening dan menunjukkan evidence hasil belajar. & Bentuk penilaian: Filtering Spasial: Konvolusi, Gaussian Blur, Sharpening. Mengacu rubrik RTM. & Ketepatan konsep; kualitas implementasi; validasi dan dokumentasi. & 3 \\
4 & SCPMK0705-04301 & Deteksi tepi: Sobel dan Canny. & Modalitas: Blended Learning. Bentuk: Luring/Daring. Metode: Case Method, praktik implementasi, eksplorasi dataset, dan peer review. & 1 x 3 x 50' tatap muka; 1 x 3 x 50' tugas/praktik mandiri. & Mahasiswa mengerjakan aktivitas terarah untuk deteksi tepi: Sobel dan Canny dan menunjukkan evidence hasil belajar. & Bentuk penilaian: Deteksi Tepi: Sobel dan Canny. Mengacu rubrik RTM. & Ketepatan konsep; kualitas implementasi; validasi dan dokumentasi. & 3 \\
5 & SCPMK0705-04301 & Kuis 1: pengolahan citra dasar. & Modalitas: Blended Learning. Bentuk: Luring/Daring. Metode: Case Method, praktik implementasi, eksplorasi dataset, dan peer review. & 1 x 3 x 50' tatap muka; 1 x 3 x 50' tugas/praktik mandiri. & Mahasiswa mengerjakan aktivitas terarah untuk kuis 1: pengolahan citra dasar dan menunjukkan evidence hasil belajar. & Bentuk penilaian: Kuis 1: Pengolahan Citra Dasar. Mengacu rubrik RTM. & Ketepatan konsep; kualitas implementasi; validasi dan dokumentasi. & 5 \\
6 & SCPMK0705-04301 & Segmentasi citra: thresholding dan watershed. & Modalitas: Blended Learning. Bentuk: Luring/Daring. Metode: Case Method, praktik implementasi, eksplorasi dataset, dan peer review. & 1 x 3 x 50' tatap muka; 1 x 3 x 50' tugas/praktik mandiri. & Mahasiswa mengerjakan aktivitas terarah untuk segmentasi citra: thresholding dan watershed dan menunjukkan evidence hasil belajar. & Bentuk penilaian: Segmentasi Citra: Thresholding dan Watershed. Mengacu rubrik RTM. & Ketepatan konsep; kualitas implementasi; validasi dan dokumentasi. & 3 \\
7 & SCPMK1004-04303 & Transformasi morfologi dan analisis bentuk. & Modalitas: Blended Learning. Bentuk: Luring/Daring. Metode: Case Method, praktik implementasi, eksplorasi dataset, dan peer review. & 1 x 3 x 50' tatap muka; 1 x 3 x 50' tugas/praktik mandiri. & Mahasiswa mengerjakan aktivitas terarah untuk transformasi morfologi dan analisis bentuk dan menunjukkan evidence hasil belajar. & Bentuk penilaian: Transformasi Morfologi dan Analisis Bentuk. Mengacu rubrik RTM. & Ketepatan konsep; kualitas implementasi; validasi dan dokumentasi. & 3 \\
8 & SCPMK0705-04301, SCPMK1004-04303 & UTS: pengolahan citra dan evaluasi kualitas. & Modalitas: Blended Learning. Bentuk: Luring/Daring. Metode: Case Method, praktik implementasi, eksplorasi dataset, dan peer review. & 1 x 3 x 50' tatap muka; 1 x 3 x 50' tugas/praktik mandiri. & Mahasiswa mengerjakan aktivitas terarah untuk UTS: pengolahan citra dan evaluasi kualitas dan menunjukkan evidence hasil belajar. & Bentuk penilaian: UTS: Pengolahan Citra dan Evaluasi Kualitas. Mengacu rubrik RTM. & Ketepatan konsep; kualitas implementasi; validasi dan dokumentasi. & 15 \\
9 & SCPMK0706-04302 & CNN: arsitektur dan transfer learning. & Modalitas: Blended Learning. Bentuk: Luring/Daring. Metode: Case Method, praktik implementasi, eksplorasi dataset, dan peer review. & 1 x 3 x 50' tatap muka; 1 x 3 x 50' tugas/praktik mandiri. & Mahasiswa mengerjakan aktivitas terarah untuk CNN: arsitektur dan transfer learning dan menunjukkan evidence hasil belajar. & Bentuk penilaian: CNN: Arsitektur dan Transfer Learning. Mengacu rubrik RTM. & Ketepatan konsep; kualitas implementasi; validasi dan dokumentasi. & 3 \\
10 & SCPMK0706-04302 & Image classification dengan CNN. & Modalitas: Blended Learning. Bentuk: Luring/Daring. Metode: Case Method, praktik implementasi, eksplorasi dataset, dan peer review. & 1 x 3 x 50' tatap muka; 1 x 3 x 50' tugas/praktik mandiri. & Mahasiswa mengerjakan aktivitas terarah untuk image classification dengan CNN dan menunjukkan evidence hasil belajar. & Bentuk penilaian: Image Classification dengan CNN. Mengacu rubrik RTM. & Ketepatan konsep; kualitas implementasi; validasi dan dokumentasi. & 3 \\
11 & SCPMK0706-04302 & Object detection: YOLO dan Faster R-CNN. & Modalitas: Blended Learning. Bentuk: Luring/Daring. Metode: Case Method, praktik implementasi, eksplorasi dataset, dan peer review. & 1 x 3 x 50' tatap muka; 1 x 3 x 50' tugas/praktik mandiri. & Mahasiswa mengerjakan aktivitas terarah untuk object detection: YOLO dan Faster R-CNN dan menunjukkan evidence hasil belajar. & Bentuk penilaian: Object Detection: YOLO dan Faster R-CNN. Mengacu rubrik RTM. & Ketepatan konsep; kualitas implementasi; validasi dan dokumentasi. & 3 \\
12 & SCPMK0706-04302 & Image segmentation: FCN dan U-Net. & Modalitas: Blended Learning. Bentuk: Luring/Daring. Metode: Case Method, praktik implementasi, eksplorasi dataset, dan peer review. & 1 x 3 x 50' tatap muka; 1 x 3 x 50' tugas/praktik mandiri. & Mahasiswa mengerjakan aktivitas terarah untuk image segmentation: FCN dan U-Net dan menunjukkan evidence hasil belajar. & Bentuk penilaian: Image Segmentation: FCN dan U-Net. Mengacu rubrik RTM. & Ketepatan konsep; kualitas implementasi; validasi dan dokumentasi. & 3 \\
13 & SCPMK1005-04304 & Kuis 2: evaluasi model visi komputer. & Modalitas: Blended Learning. Bentuk: Luring/Daring. Metode: Case Method, praktik implementasi, eksplorasi dataset, dan peer review. & 1 x 3 x 50' tatap muka; 1 x 3 x 50' tugas/praktik mandiri. & Mahasiswa mengerjakan aktivitas terarah untuk kuis 2: evaluasi model visi komputer dan menunjukkan evidence hasil belajar. & Bentuk penilaian: Kuis 2: Evaluasi Model Visi Komputer. Mengacu rubrik RTM. & Ketepatan konsep; kualitas implementasi; validasi dan dokumentasi. & 5 \\
14 & SCPMK0706-04302 & Data augmentation dan transfer learning lanjutan. & Modalitas: Blended Learning. Bentuk: Luring/Daring. Metode: Case Method, praktik implementasi, eksplorasi dataset, dan peer review. & 1 x 3 x 50' tatap muka; 1 x 3 x 50' tugas/praktik mandiri. & Mahasiswa mengerjakan aktivitas terarah untuk data augmentation dan transfer learning lanjutan dan menunjukkan evidence hasil belajar. & Bentuk penilaian: Data Augmentation dan Transfer Learning Lanjutan. Mengacu rubrik RTM. & Ketepatan konsep; kualitas implementasi; validasi dan dokumentasi. & 3 \\
15 & SCPMK1005-04304 & Studi kasus: aplikasi visi komputer industri. & Modalitas: Blended Learning. Bentuk: Luring/Daring. Metode: Case Method, praktik implementasi, eksplorasi dataset, dan peer review. & 1 x 3 x 50' tatap muka; 1 x 3 x 50' tugas/praktik mandiri. & Mahasiswa mengerjakan aktivitas terarah untuk studi kasus: aplikasi visi komputer industri dan menunjukkan evidence hasil belajar. & Bentuk penilaian: Studi Kasus: Aplikasi Visi Komputer Industri. Mengacu rubrik RTM. & Ketepatan konsep; kualitas implementasi; validasi dan dokumentasi. & 3 \\
16 & SCPMK0706-04302, SCPMK1005-04304 & PBL: implementasi sistem visi komputer. & Modalitas: Blended Learning. Bentuk: Luring/Daring. Metode: Case Method, praktik implementasi, eksplorasi dataset, dan peer review. & 1 x 3 x 50' tatap muka; 1 x 3 x 50' tugas/praktik mandiri. & Mahasiswa mengerjakan aktivitas terarah untuk PBL: implementasi sistem visi komputer dan menunjukkan evidence hasil belajar. & Bentuk penilaian: PBL: Implementasi Sistem Visi Komputer. Mengacu rubrik RTM. & Ketepatan konsep; kualitas implementasi; validasi dan dokumentasi. & 25 \\
17 & SCPMK0706-04302, SCPMK1005-04304 & UAS: presentasi dan evaluasi sistem visi komputer. & Modalitas: Blended Learning. Bentuk: Luring/Daring. Metode: Case Method, praktik implementasi, eksplorasi dataset, dan peer review. & 1 x 3 x 50' tatap muka; 1 x 3 x 50' tugas/praktik mandiri. & Mahasiswa mengerjakan aktivitas terarah untuk UAS: presentasi dan evaluasi sistem visi komputer dan menunjukkan evidence hasil belajar. & Bentuk penilaian: UAS: Presentasi dan Evaluasi Sistem Visi Komputer. Mengacu rubrik RTM. & Ketepatan konsep; kualitas implementasi; validasi dan dokumentasi. & 22 \\
\end{longtblr}
